% LAMPIRAN
% Setelah \appendix, chapter akan otomatis menggunakan huruf (A, B, C, dst)

\chapter{Kode Program Python untuk Analisis Spektrum}
\label{app:kode-python}
\appendices{Kode Program Python untuk Analisis Spektrum}

Berikut adalah contoh kode program Python yang digunakan untuk menganalisis spektrum plasma.

\begin{verbatim}
import numpy as np
import matplotlib.pyplot as plt

def analyze_spectrum(wavelength, intensity):
    """
    Fungsi untuk menganalisis spektrum
    """
    peaks = find_peaks(intensity)
    return peaks

# Load data
data = np.loadtxt('spectrum_data.txt')
wavelength = data[:, 0]
intensity = data[:, 1]

# Analisis
peaks = analyze_spectrum(wavelength, intensity)
print(f"Ditemukan {len(peaks)} puncak spektrum")

# Plot hasil
plt.figure(figsize=(10, 6))
plt.plot(wavelength, intensity)
plt.xlabel('Panjang Gelombang (nm)')
plt.ylabel('Intensitas (a.u.)')
plt.title('Spektrum Plasma Hidrogen')
plt.grid(True)
plt.show()
\end{verbatim}

\section{Penjelasan Kode}

Program di atas melakukan analisis spektrum dengan langkah-langkah berikut:
\begin{enumerate}
    \item Membaca data spektrum dari file
    \item Mencari puncak-puncak intensitas
    \item Memvisualisasikan hasil dalam bentuk grafik
\end{enumerate}

\blindtext[2]

\chapter{Data Hasil Pengukuran}
\label{app:data-pengukuran}
\appendices{Data Hasil Pengukuran}

Berikut adalah tabel data hasil pengukuran spektrum plasma pada berbagai kondisi.

\section{Data Temperatur Plasma}

\begin{table}[h]
\centering
\caption{Hasil pengukuran temperatur plasma}
\begin{tabular}{|c|c|c|c|}
\hline
\textbf{No.} & \textbf{Tekanan (Torr)} & \textbf{Daya (W)} & \textbf{Temperatur (K)} \\ \hline
1 & 0.5 & 100 & 8500 \\ \hline
2 & 0.5 & 150 & 9200 \\ \hline
3 & 0.5 & 200 & 10100 \\ \hline
4 & 1.0 & 100 & 7800 \\ \hline
5 & 1.0 & 150 & 8600 \\ \hline
6 & 1.0 & 200 & 9500 \\ \hline
7 & 1.5 & 100 & 7200 \\ \hline
8 & 1.5 & 150 & 8100 \\ \hline
9 & 1.5 & 200 & 9000 \\ \hline
\end{tabular}
\end{table}

\blindtext[2]

\section{Data Kerapatan Elektron}

\begin{table}[h]
\centering
\caption{Hasil pengukuran kerapatan elektron}
\begin{tabular}{|c|c|c|c|}
\hline
\textbf{No.} & \textbf{Tekanan (Torr)} & \textbf{Daya (W)} & \textbf{$n_e$ ($\times 10^{16}$ cm$^{-3}$)} \\ \hline
1 & 0.5 & 100 & 2.3 \\ \hline
2 & 0.5 & 150 & 3.1 \\ \hline
3 & 0.5 & 200 & 4.2 \\ \hline
4 & 1.0 & 100 & 3.5 \\ \hline
5 & 1.0 & 150 & 4.8 \\ \hline
6 & 1.0 & 200 & 6.1 \\ \hline
7 & 1.5 & 100 & 4.2 \\ \hline
8 & 1.5 & 150 & 5.6 \\ \hline
9 & 1.5 & 200 & 7.3 \\ \hline
\end{tabular}
\end{table}

\blindtext

\chapter{Dokumentasi Peralatan}
\label{app:dokumentasi}
\appendices{Dokumentasi Peralatan}

\section{Spesifikasi Spektrometer}

Spektrometer yang digunakan dalam penelitian ini memiliki spesifikasi sebagai berikut:

\begin{itemize}
    \item \textbf{Model}: Ocean Optics HR4000
    \item \textbf{Range panjang gelombang}: 200-1100 nm
    \item \textbf{Resolusi optik}: 0.025 nm (FWHM)
    \item \textbf{Detector}: Sony ILX511B CCD array (3648 pixels)
    \item \textbf{Signal-to-noise ratio}: 250:1 (single acquisition)
    \item \textbf{Integration time}: 1 ms - 65 seconds
\end{itemize}

\section{Diagram Setup Eksperimen}

\blindtext[2]

% Jika ada gambar, bisa ditambahkan seperti ini:
% \begin{figure}[h]
% \centering
% \includegraphics[width=0.8\textwidth]{images/setup-eksperimen.pdf}
% \caption{Diagram setup eksperimen}
% \label{fig:setup-eksperimen}
% \end{figure}

\section{Prosedur Kalibrasi}

\begin{enumerate}
    \item Nyalakan spektrometer dan tunggu 30 menit untuk warm-up
    \item Lakukan dark spectrum measurement
    \item Gunakan lampu kalibrasi Mercury-Argon
    \item Identifikasi minimal 5 garis kalibrasi
    \item Lakukan polynomial fitting untuk kalibrasi panjang gelombang
    \item Simpan file kalibrasi untuk digunakan pada pengukuran selanjutnya
\end{enumerate}

\blindtext[3]
